\begin{tikzpicture}[node distance=2cm, transform shape, on grid, ->]
    \tikzstyle{every node}=[font=\small]
    \node (start) [code] {\lstinline|beschreibe_wetter(19)|};
    \node (dec1) [decision, below of=start] {\lstinline|temperatur >= 30|};
    \node (printif1) [code, text width=4cm, right=1cm of dec1.east] {\lstinline|print("Es ist heiss")|};
    \node (dec2) [decision, below=2.5cm of dec1] {\lstinline|temperatur >= 20|};
    \node (printif2) [code, text width=4cm, right=1cm of dec2.east] {\lstinline|print("Es ist warm")|};
    \node (printelse) [code, below of=dec2] {\lstinline|print("Es ist kühl")|};
    \node (end) [circle,draw=black, fill=black, below of=printelse, yshift=.8cm] {};
    \node (intermediate1) [coordinate, below=1.5cm of end] {};
    \node (intermediate2) [coordinate, left=3.5cm of dec2] {};
    
    \draw (start) -- (dec1);
    \draw (dec1) -- node[anchor=top, above=.1cm] {\lstinline|True|} (printif1);
    \draw (dec1) -- node[anchor=east] {\lstinline|False|} (dec2);
    
    \draw (dec2) -- node[anchor=top, above=.1cm] {\lstinline|True|} (printif2);
    \draw (dec2) -- node[anchor=east] {\lstinline|False|} (printelse);
    
    \draw[red, thick] (printif1) --++ (3.5cm,0cm) |- (intermediate1) -| (intermediate2) -- (dec2);
    \draw (printif2) --++ (2.5cm,0cm) |- (end);
    \draw (printelse) -- (end);
    
    \draw (printelse) -- (end);
\end{tikzpicture}